\documentclass[a4paper,12pt]{article}
\usepackage[utf8]{inputenc}
\usepackage[T1]{fontenc}
\usepackage[ngerman]{babel}
\usepackage{amsmath, amssymb}
\usepackage{geometry}
\usepackage{parskip}
\usepackage{titlesec}
\usepackage{enumitem}
\usepackage{array}
\usepackage{booktabs}

\geometry{left=3cm, right=3cm, top=2.5cm, bottom=2.5cm}

\titleformat{\section}{\large\bfseries}{}{0em}{}[\titlerule]
\titleformat{\subsection}{\normalsize\bfseries\scshape}{}{0em}{}

% Glossar-Eintrag: Nummer, fetter Titel, Inhalt
\newcommand{\gentry}[3]{%
  \noindent\textbf{#1\quad #2}\nopagebreak\\[0.2em]
  #3\par\smallskip
}

\begin{document}

\begin{center}
  {\LARGE\bfseries Glossar zu Lektion 2}\\[0.5em]
  {\large Mathematische Grundlagen (Modul 61111)}
\end{center}

\vspace{1em}

Dieses Glossar fasst die wesentlichen Definitionen, Merkregeln und Sätze der
zweiten Lektion kompakt zusammen. Nummerierungen entsprechen dem Lehrtext.

% ======================================================================
\section*{4.1 \quad Treppennormalformen}

\gentry{4.1.1}{Treppennormalform}{%
  Eine Matrix $A \in M_{mn}(K)$ ist in \textbf{Treppennormalform}, wenn $A$ die
  Nullmatrix ist oder wenn es $r$ Spaltenindizes $j_1 < j_2 < \cdots < j_r$ gibt,
  sodass für alle $1 \le i \le r$ gilt:
  \begin{enumerate}[label=(\roman*),topsep=2pt,itemsep=1pt]
    \item $a_{ij_i} = 1$ (Pivot-Eintrag),
    \item $a_{lj_i} = 0$ für alle $l \neq i$ (Nullen in der Pivot-Spalte),
    \item $a_{il} = 0$ für alle $l < j_i$ (Nullen links des Pivots),
    \item $a_{kl} = 0$ für alle $k > r$ und alle $1 \le l \le n$ (Nullzeilen am Ende).
  \end{enumerate}
}

\gentry{4.1.2}{Pivot-Positionen}{%
  Die Paare $(1,j_1),(2,j_2),\ldots,(r,j_r)$ der Treppennormalform heißen
  \textbf{Pivot-Positionen} (gesprochen: Pivo-Positionen).
}

% ======================================================================
\section*{4.2 \quad Der Gaußalgorithmus}

\gentry{4.2}{Gaußalgorithmus}{%
  Systematisches Verfahren, das jede Matrix $A \in M_{mn}(K)$ durch endlich
  viele elementare Zeilenumformungen in Treppennormalform überführt.
  Man arbeitet spaltenweise von links nach rechts:
  \begin{enumerate}[label=\arabic*.,topsep=2pt,itemsep=1pt]
    \item Falls alle Einträge $a_{it} = 0$ für $i \ge r$: Spalte überspringen.
    \item Falls $a_{rt} \neq 0$: Zeile $r$ auf $1$ normieren, dann alle anderen
          Einträge in Spalte $t$ auf $0$ bringen.
    \item Falls $a_{rt} = 0$, aber $a_{it} \neq 0$ für ein $i > r$:
          Zeile $r$ und Zeile $i$ tauschen, dann wie in Schritt~2 vorgehen.
  \end{enumerate}
}

\gentry{4.2.4}{Satz (Existenz von Treppennormalformen)}{%
  Sei $A \in M_{mn}(K)$. Dann gibt es Elementarmatrizen $E_1,\ldots,E_s \in M_{mm}(K)$,
  so dass $E_s E_{s-1} \cdots E_1 A$ eine Matrix in Treppennormalform ist.
}

% ======================================================================
\section*{4.3 \quad Transformationsmatrix}

\gentry{4.3}{Transformationsmatrix $S$}{%
  Zur Berechnung von $S$ mit $SA = T$ (Treppennormalform) schreibt man
  $\bar{A} = \bigl(A \mid I_m\bigr)$ und führt auf $\bar{A}$ den Gaußalgorithmus durch.
  Das Ergebnis ist $\bar{A}' = \bigl(T \mid S\bigr)$.
  \smallskip\\
  $S$ ist invertierbar (Produkt von Elementarmatrizen).
}

% ======================================================================
\section*{4.4 \quad Eindeutigkeit der Treppennormalform}

\gentry{4.4.1}{Lemma (Invertierbar mit $GT = T'$ in TNF)}{%
  Seien $T, T' \in M_{mn}(K)$ in Treppennormalform, und sei $G \in M_{mm}(K)$
  invertierbar mit $GT = T'$. Dann gilt $T = T'$.
}

\gentry{4.4.2}{Satz (Eindeutigkeit der Treppennormalform)}{%
  Sei $A \in M_{mn}(K)$. Alle Matrizen in Treppennormalform, die durch elementare
  Zeilenumformungen aus $A$ entstehen, sind gleich.
  \smallskip\\
  Man nennt diese eindeutige Matrix die \textbf{Treppennormalform zu $A$}.
}

\gentry{4.4.5}{Korollar (Zeilenäquivalenz und Treppennormalform)}{%
  Seien $A, B \in M_{mn}(K)$. Dann gilt:
  $A$ und $B$ sind genau dann zeilenäquivalent, wenn sie dieselbe Treppennormalform haben.
}

% ======================================================================
\section*{4.5 \quad Rang einer Matrix}

\gentry{4.5.1}{Rang}{%
  Sei $T$ die Treppennormalform zu $A \in M_{mn}(K)$. Die Anzahl der Pivot-Positionen
  in $T$ heißt \textbf{Rang} von $A$: $\operatorname{Rg}(A)$.
}

\gentry{4.5.4}{Proposition (Invertierbare Matrizen)}{%
  Sei $A \in M_{mm}(K)$. Folgende Aussagen sind äquivalent:
  \begin{enumerate}[label=(\alph*),topsep=2pt,itemsep=1pt]
    \item $A$ ist invertierbar.
    \item Die Treppennormalform zu $A$ ist $I_m$.
    \item $\operatorname{Rg}(A) = m$.
    \item $A$ ist ein Produkt von Elementarmatrizen.
  \end{enumerate}
}

\gentry{4.5.5}{Proposition (Ränge von Matrizen)}{%
  Sei $A \in M_{mn}(K)$. Dann gilt:
  \begin{itemize}[topsep=2pt,itemsep=1pt]
    \item Ist $P \in M_{mm}(K)$ invertierbar, so gilt $\operatorname{Rg}(A) = \operatorname{Rg}(PA)$.
    \item Ist $Q \in M_{nn}(K)$ invertierbar, so gilt $\operatorname{Rg}(A) = \operatorname{Rg}(AQ)$.
    \item Ist $A = A'A''$, so gilt $\operatorname{Rg}(A) \le \operatorname{Rg}(A')$.
  \end{itemize}
}

\gentry{4.5.6}{Korollar (Linksinverse Matrix ist Inverse)}{%
  Seien $A, C \in M_{mm}(K)$ mit $CA = I_m$.
  Dann sind $C$ und $A$ invertierbar, und es gilt $C = A^{-1}$.
}

% ======================================================================
\section*{5.1 \quad Lineare Gleichungssysteme}

\gentry{5.1.1}{Lineares Gleichungssystem}{%
  Ein \textbf{lineares Gleichungssystem} über $K$ mit $m$ Zeilen und $n$
  Unbestimmten $x_1,\ldots,x_n$ hat die Form $Ax = b$, wobei
  $A = (a_{ij}) \in M_{mn}(K)$ die \textbf{Koeffizientenmatrix} ist.
}

\gentry{5.1.4}{Erweiterte Koeffizientenmatrix}{%
  Die $m \times (n+1)$-Matrix $A' = \bigl(A \mid b\bigr)$ heißt
  \textbf{erweiterte Koeffizientenmatrix} von $Ax = b$.
  \smallskip\\
  Das Gleichungssystem heißt \textbf{homogen}, falls $b = 0$, sonst \textbf{inhomogen}.
}

% ======================================================================
\section*{5.2 \quad Struktur der Lösungsmenge}

\gentry{5.2.1}{Proposition (Lösungsmenge als $\lambda_0 + U$)}{%
  Sei $\lambda_0$ eine Lösung von $Ax = b$, und sei $U$ die Lösungsmenge des
  zugehörigen homogenen Systems $Ax = 0$. Dann gilt:
  \[
    L = \lambda_0 + U := \{\lambda_0 + u \mid u \in U\}.
  \]
}

\gentry{5.2.3}{Proposition (Lösungsmenge invariant unter $P$)}{%
  Sei $P \in M_{mm}(K)$ invertierbar. Die Lösungsmengen von $Ax = b$ und $(PA)x = Pb$ stimmen überein.
}

\gentry{5.2.4}{Korollar (Zeilenäquivalente Systeme, gleiche Lösungen)}{%
  Sind $A'$ und $B'$ zeilenäquivalente erweiterte Koeffizientenmatrizen, so
  haben die zugehörigen Gleichungssysteme dieselbe Lösungsmenge.
}

\gentry{5.2.5}{Proposition (Lösbarkeit)}{%
  Das System $Ax = b$ hat genau dann mindestens eine Lösung, wenn
  \[
    \operatorname{Rg}(A) = \operatorname{Rg}(A').
  \]
}

\gentry{5.2.9}{Merkregel (partikuläre Lösung)}{%
  Sei $T$ in Treppennormalform, $\operatorname{Rg}(T) = \operatorname{Rg}(T')$.
  \begin{enumerate}[label=\arabic*.,topsep=2pt,itemsep=1pt]
    \item Streiche alle Nullzeilen von $T'$.
    \item Füge Nullzeilen so ein, dass die Matrix links des Striches quadratisch
          wird und die Pivot-Elemente Diagonalelemente werden.
  \end{enumerate}
  Dann steht rechts des Striches eine partikuläre Lösung $\lambda_0$.
}

\gentry{5.2.11}{Satz (Struktur der Lösungsmenge)}{%
  Sei $U$ die Lösungsmenge von $Ax = 0$, und sei $L$ die Lösungsmenge von $Ax = b$.
  \begin{enumerate}[label=(\alph*),topsep=2pt,itemsep=1pt]
    \item $Ax = b$ ist lösbar $\iff$ $\operatorname{Rg}(A) = \operatorname{Rg}(A')$.
    \item Ist $\lambda_0$ eine Lösung, so gilt $L = \lambda_0 + U$.
    \item $L$ und die Lösungsmenge von $Tx = d$ (Treppennormalform) stimmen überein.
  \end{enumerate}
}

\gentry{5.2.14}{Korollar (Eindeutige Lösung)}{%
  Besitzt $Ax = b$ eine Lösung, so ist sie genau dann eindeutig, wenn
  $\operatorname{Rg}(A) = n$ (Anzahl der Unbestimmten).
}

\gentry{5.2.16}{Merkregel (homogene Lösungsmenge)}{%
  Sei $T$ in Treppennormalform, $r = \operatorname{Rg}(T) < n$.
  \begin{enumerate}[label=\arabic*.,topsep=2pt,itemsep=1pt]
    \item Streiche alle Nullzeilen von $T$.
    \item Füge Nullzeilen ein, sodass die Matrix quadratisch wird und Pivot-Elemente
          auf der Diagonalen stehen.
    \item Ersetze die Nullen auf der Diagonalen durch $-1$.
    \item Die Spalten $S_1,\ldots,S_{n-r}$, in die $-1$ eingefügt wurde, spannen $U$ auf:
          \[U = \{a_1 S_1 + \cdots + a_{n-r} S_{n-r} \mid a_i \in K\}.\]
  \end{enumerate}
  Ist $r = n$, so gilt $U = \{0\}$.
}

% ======================================================================
\section*{6.1 \quad Vektorräume}

\gentry{6.1.1}{Vektorraum}{%
  Ein \textbf{$K$-Vektorraum} $V$ ist eine Menge mit einer Addition
  $+: V \times V \to V$ und einer Skalarmultiplikation $\cdot: K \times V \to V$,
  sodass folgende Axiome gelten:
  \begin{enumerate}[label=(\roman*),topsep=2pt,itemsep=1pt]
    \item Addition: kommutativ, assoziativ; neutrales Element $0 \in V$;
          jedes $v \in V$ hat ein inverses Element $-v$.
    \item Skalarmultiplikation: $(ab)v = a(bv)$;\quad $1 \cdot v = v$.
    \item Distributivgesetze: $a(v_1+v_2) = av_1+av_2$;\quad $(a+b)v = av+bv$.
  \end{enumerate}
  Elemente von $V$ heißen \textbf{Vektoren}, Elemente von $K$ heißen
  \textbf{Skalare}.
}

\gentry{6.1.4}{Proposition (Rechenregeln)}{%
  Sei $V$ ein $K$-Vektorraum. Dann gilt:
  \begin{enumerate}[label=(\alph*),topsep=2pt,itemsep=1pt]
    \item $0 \cdot v = 0$ für alle $v \in V$.
    \item $a \cdot 0 = 0$ für alle $a \in K$.
    \item $(-a)v = -(av) = a(-v)$ für alle $a \in K$, $v \in V$.
  \end{enumerate}
}

\gentry{}{Beispiele für Vektorräume}{%
  \begin{itemize}[topsep=2pt,itemsep=1pt]
    \item $M_{mn}(K)$ mit Matrix-Addition und Skalarmultiplikation.
    \item $K^n = M_{n1}(K)$: Spalten der Länge $n$.
    \item Lösungsmenge $U$ eines homogenen Systems $Ax = 0$ (Unterraum von $K^n$).
    \item $K[T]$: Polynome über $K$ (nicht endlich erzeugt).
    \item Polynome vom Grad $\le n$: endlich erzeugt.
  \end{itemize}
}

% ======================================================================
\section*{6.2 \quad Unterräume}

\gentry{6.2.1}{Unterraum}{%
  Eine nicht leere Teilmenge $U \neq \emptyset$ eines $K$-Vektorraums $V$ heißt
  \textbf{Unterraum}, wenn $U$ mit den Verknüpfungen aus $V$ selbst ein Vektorraum ist.
}

\gentry{6.2.3}{Proposition (Unterraumkriterium)}{%
  $U \subseteq V$ ist genau dann ein Unterraum, wenn:
  \begin{enumerate}[label=(\alph*),topsep=2pt,itemsep=1pt]
    \item $0 \in U$,
    \item $u_1 + u_2 \in U$ für alle $u_1, u_2 \in U$,
    \item $au \in U$ für alle $a \in K$, $u \in U$.
  \end{enumerate}
}

\gentry{6.2.7}{Proposition (Durchschnitt)}{%
  Sind $U$ und $W$ Unterräume von $V$, so ist $U \cap W$ ein Unterraum von $V$.
}

\gentry{6.2.9}{Summe von Unterräumen}{%
  Sind $U$ und $W$ Unterräume von $V$, so ist
  \[
    U + W := \{u + w \mid u \in U,\, w \in W\}
  \]
  ein Unterraum von $V$. Es gilt $U \cup W \subseteq U + W$.
}

% ======================================================================
\section*{6.3 \quad Endlich erzeugte Vektorräume}

\gentry{6.3.1}{Linearkombination}{%
  $a_1 v_1 + \cdots + a_m v_m$ heißt \textbf{Linearkombination} der Vektoren
  $v_1,\ldots,v_m \in V$ mit Koeffizienten $a_i \in K$.
}

\gentry{6.3.2}{Lineare Hülle (Erzeugnis)}{%
  Die Menge aller Linearkombinationen endlich vieler Vektoren aus $S \subseteq V$
  heißt \textbf{lineare Hülle} oder \textbf{Erzeugnis} von $S$:
  \[
    \langle S \rangle. \qquad \text{Für } S = \emptyset \text{ setzt man } \langle \emptyset \rangle := \{0\}.
  \]
  $\langle S \rangle$ ist stets ein Unterraum von $V$ (Prop. 6.3.4).
}

\gentry{6.3.6}{Erzeugendensystem}{%
  $S \subseteq V$ heißt \textbf{Erzeugendensystem} von $V$, wenn $\langle S \rangle = V$.
}

\gentry{6.3.8}{Endlich erzeugter Vektorraum}{%
  $V$ heißt \textbf{endlich erzeugt}, wenn es eine endliche Menge
  $\{v_1,\ldots,v_m\}$ mit $V = \langle v_1,\ldots,v_m \rangle$ gibt.
  \smallskip\\
  Beispiele: $K^n$, $M_{mn}(K)$, Polynome vom Grad $\le n$, Lösungsmenge von $Ax=0$.
  \smallskip\\
  Gegenbeispiel: $K[T]$ ist \emph{nicht} endlich erzeugt.
}

\end{document}
